% !TEX root = ../algebra_02.tex

\section{Равенство столбцового и строчного ранга}

\begin{Theorem}{Равенство рангов}{}
	Для любой матрицы $A \in k^{m \times n}$: \qquad $\mathrm{rk}_{\mathrm{col}}(A) = \mathrm{rk}_{\mathrm{row}}(A).$
	
	Общую величину называют просто \emph{рангом} матрицы и обозначают $\mathrm{rk}(A)$.
\end{Theorem}

\begin{proof}
	Элементарными преобразованиями строк и столбцов приведём $A$ к нормальному виду:
	\[
	D = \begin{pmatrix} E_r & 0 \\ 0 & 0 \end{pmatrix},
	\]
	где $r$ — число ненулевых строк. 
	
	Как было доказано ранее, элементарные преобразования строк не меняют ни столбцовый, ни строковый ранг. Аналогично, элементарные преобразования столбцов также сохраняют оба ранга: они не меняют линейную оболочку столбцов и обратимо преобразуют столбцы.
	
	Значит, ранги матрицы $A$ совпадают с соответствующими рангами матрицы $D$.
	Для матрицы $D$ утверждение теоремы очевидно:
	\begin{itemize}
		\item $\mathrm{rk}_{\mathrm{col}}(D) = r$ (ровно $r$ ненулевых столбцов, каждый из которых содержит единственную единицу, поэтому они линейно независимы).
		\item $\mathrm{rk}_{\mathrm{row}}(D) = r$ (ровно $r$ ненулевых строк, которые также образуют стандартный базис, поэтому они линейно независимы).
	\end{itemize}
	Следовательно, $\mathrm{rk}_{\mathrm{col}}(A) = r = \mathrm{rk}_{\mathrm{row}}(A)$.
\end{proof}
